% =============================================================================
%  week15_diffusion.tex
%  Diffusion Models: From VAEs to DDPM — Full Mathematical Derivation
%  Morten Hjorth-Jensen, May 2026
%  Revised: code slides removed; VAE kept as motivation; new derivations added
% =============================================================================

\documentclass[aspectratio=169,11pt]{beamer}

\usetheme{Madrid}
\usecolortheme{default}
\usefonttheme[onlymath]{serif}

\setbeamercolor{frametitle}{bg=blue!11!white, fg=blue!60!black}
\setbeamercolor{title}{fg=blue!65!black}
\setbeamercolor{block title}{bg=blue!16!white, fg=blue!60!black}
\setbeamercolor{block body}{bg=blue!4!white}
\setbeamercolor{alerted text}{fg=red!65!black}
\setbeamertemplate{navigation symbols}{}
\setbeamertemplate{footline}{%
  \leavevmode\hbox{%
    \begin{beamercolorbox}[wd=.30\paperwidth,ht=2.25ex,dp=1ex,center]{author in head/foot}%
      \usebeamerfont{author in head/foot}\insertshortauthor
    \end{beamercolorbox}%
    \begin{beamercolorbox}[wd=.44\paperwidth,ht=2.25ex,dp=1ex,center]{title in head/foot}%
      \usebeamerfont{title in head/foot}\insertshorttitle
    \end{beamercolorbox}%
    \begin{beamercolorbox}[wd=.26\paperwidth,ht=2.25ex,dp=1ex,right]{date in head/foot}%
      \usebeamerfont{date in head/foot}\insertshortdate{}\hspace*{1em}%
      \insertframenumber/\inserttotalframenumber\hspace*{1ex}
    \end{beamercolorbox}}%
  \vskip0pt}

\usepackage[utf8]{inputenc}
\usepackage[T1]{fontenc}
\usepackage{amsmath,amssymb,amsfonts,mathtools,bm}
\usepackage{booktabs}
\usepackage{xcolor}
\usepackage{tikz}
\usetikzlibrary{arrows.meta,positioning,shapes.geometric,decorations.pathreplacing}
\usepackage{hyperref}

% ── Convenience macros ───────────────────────────────────────────────────────
\newcommand{\E}{\mathbb{E}}
\newcommand{\R}{\mathbb{R}}
\newcommand{\DKL}{D_{\mathrm{KL}}}
\newcommand{\Norm}{\mathcal{N}}
\newcommand{\ELBO}{\mathcal{L}}
\newcommand{\bx}{\bm{x}}
\newcommand{\bh}{\bm{h}}
\newcommand{\bz}{\bm{z}}
\newcommand{\beps}{\bm{\epsilon}}
\newcommand{\bmu}{\bm{\mu}}
\newcommand{\bsig}{\bm{\sigma}}
\newcommand{\btheta}{\bm{\theta}}
\newcommand{\bphi}{\bm{\phi}}
\newcommand{\half}{\tfrac{1}{2}}
\DeclareMathOperator*{\argmin}{arg\,min}

% ── Metadata ─────────────────────────────────────────────────────────────────
\title[Diffusion Models]{Denoising Diffusion Probabilistic Models:\\
       A Complete Mathematical Derivation}
\subtitle{From Variational Autoencoders to Score-Based Generative Models}
\author[M.\ Hjorth-Jensen]{Morten Hjorth-Jensen}
\institute[UiO/MSU]{Department of Physics \& CCSE, University of Oslo\\
                    \& Michigan State University}
\date{May 2026}

% =============================================================================
\begin{document}
% =============================================================================

\begin{frame}[plain]
  \titlepage
\end{frame}

% ── Table of contents ─────────────────────────────────────────────────────────
\begin{frame}{Outline}
  \tableofcontents
\end{frame}

% =============================================================================
\section{Motivation: VAEs and Their Limits}
% =============================================================================

% -----------------------------------------------------------------------------
\begin{frame}{Generative models: the common thread}

All deep generative models share one goal: learn a distribution
$p_{\btheta}(\bx)$ over data~$\bx$ that we can both evaluate and sample from.

\medskip
\begin{center}
\renewcommand{\arraystretch}{1.25}
\begin{tabular}{llll}
\toprule
\textbf{Family} & \textbf{Latent} & \textbf{Training signal} & \textbf{Sample cost}\\
\midrule
VAE            & $\bh \in \R^{d_h},\;d_h\!\ll\!d_x$ & ELBO             & 1 decoder pass\\
\textbf{Diffusion} & $\bx_t\in\R^{d_x}$ (full dim)  & \textbf{ELBO, $T$ levels} & $T$ denoiser calls\\
GAN            & $\bz\in\R^k$                       & Adversarial      & 1 generator pass\\
\bottomrule
\end{tabular}
\end{center}

\medskip
\begin{alertblock}{Key insight}
Diffusion models are \textbf{hierarchical VAEs} with $T$ latent levels,
a fixed (non-learned) encoder, and full-dimensional latents.
Understanding VAEs is the fastest route to understanding diffusion.
\end{alertblock}
\end{frame}

% -----------------------------------------------------------------------------
\begin{frame}{The VAE: ELBO and reparameterisation (recap)}

\textbf{Goal.} Maximise $\log p_{\btheta}(\bx)$ — intractable because
$p(\bx)=\int p_{\btheta}(\bx|\bh)\,p(\bh)\,d\bh$.

\medskip
Introduce approximate posterior $q_{\bphi}(\bh|\bx)$.
By Jensen's inequality on $\log$:
\[
  \log p(\bx)
  = \log\E_{q_{\bphi}(\bh|\bx)}\!\left[\frac{p_{\btheta}(\bx,\bh)}{q_{\bphi}(\bh|\bx)}\right]
  \;\geq\;
  \E_{q_{\bphi}(\bh|\bx)}\!\left[\log\frac{p_{\btheta}(\bx,\bh)}{q_{\bphi}(\bh|\bx)}\right]
  =: \ELBO(\bphi,\btheta).
\]
Splitting the joint $p(\bx,\bh)=p_{\btheta}(\bx|\bh)\,p(\bh)$:
\[
\boxed{
\ELBO(\bphi,\btheta)
= \underbrace{\E_{q_{\bphi}(\bh|\bx)}\!\left[\log p_{\btheta}(\bx|\bh)\right]}_{\text{reconstruction}}
- \underbrace{\DKL\!\bigl(q_{\bphi}(\bh|\bx)\,\|\,p(\bh)\bigr)}_{\text{regularisation}}}
\]

\textbf{Reparameterisation.}  With Gaussian encoder
$q_{\bphi}(\bh|\bx)=\Norm(\bmu_{\bphi},\bsig^2_{\bphi}\mathbf{I})$
and prior $p(\bh)=\Norm(\bm{0},\mathbf{I})$:
\[
  \bh = \bmu_{\bphi}(\bx) + \bsig_{\bphi}(\bx)\odot\beps,\quad
  \beps\sim\Norm(\bm{0},\mathbf{I}).
\]
Closed-form KL:
$\DKL(q\|p) = \frac{1}{2}\sum_{j}\!\left(\mu_j^2+\sigma_j^2-\log\sigma_j^2-1\right)$.
\end{frame}

% -----------------------------------------------------------------------------
\begin{frame}{Why VAEs produce blurry samples}

\textbf{The $\ell_2$-averaging problem.}
The decoder minimises
\[
  -\E_{q_{\bphi}}\!\left[\log p_{\btheta}(\bx|\bh)\right]
  \;\propto\;
  \E_{q_{\bphi}}\!\left[\|\bx - \hat{\bx}_{\btheta}(\bh)\|^2\right].
\]
Because the posterior $q_{\bphi}(\bh|\bx)$ has nonzero variance,
$\hat{\bx}_{\btheta}$ learns the \emph{mean over all likely completions},
which is systematically blurry.

\bigskip
\begin{columns}[T]
\column{0.5\textwidth}
\begin{block}{Other VAE limitations}
\begin{itemize}
  \item \textbf{Posterior collapse}: powerful decoder can ignore $\bh$,
        KL$\to 0$, latent useless.
  \item \textbf{ELBO gap}:
        $\log p(\bx) = \ELBO + \DKL(q_{\bphi}\|p(\bh|\bx)) > \ELBO$.
  \item Mean-field Gaussian approximation may miss
        multi-modal true posteriors.
\end{itemize}
\end{block}

\column{0.47\textwidth}
\begin{alertblock}{The diffusion solution}
\begin{itemize}
  \item Replace 1 latent level with $T\gg 1$.
  \item Fix the encoder: no collapse possible.
  \item Predict only the \emph{noise added at one step},
        not the full image — eliminates $\ell_2$-averaging artefacts.
\end{itemize}
\end{alertblock}
\end{columns}
\end{frame}

% =============================================================================
\section{Hierarchical VAEs and the Diffusion Limit}
% =============================================================================

% -----------------------------------------------------------------------------
\begin{frame}{Markovian hierarchical VAE (MHVAE)}

Generalise VAE to a \textbf{chain of $T$ latent spaces}
$\bx_1,\bx_2,\ldots,\bx_T$ (we already write $\bx_t$ for the $t$-th latent):

\begin{align*}
p(\bx_0,\bx_{1:T})
  &= p(\bx_T)\,p_{\btheta}(\bx_0|\bx_1)
     \prod_{t=2}^{T}p_{\btheta}(\bx_{t-1}|\bx_t)
     \quad\text{(generative)}\\[4pt]
q(\bx_{1:T}|\bx_0)
  &= q(\bx_1|\bx_0)
     \prod_{t=2}^{T}q(\bx_t|\bx_{t-1})
     \quad\text{(inference)}
\end{align*}

Substituting into the ELBO and using the Markov property:
\[
\log p(\bx_0)
\geq
\E_{q(\bx_{1:T}|\bx_0)}
\!\left[\log\frac{
    p(\bx_T)\,p_{\btheta}(\bx_0|\bx_1)
    \prod_{t=2}^{T}p_{\btheta}(\bx_{t-1}|\bx_t)}
  {q(\bx_1|\bx_0)\prod_{t=2}^{T}
    q(\bx_t|\bx_{t-1})}\right]
\]
Each level contributes a reconstruction/regularisation pair.
\end{frame}

% -----------------------------------------------------------------------------
\begin{frame}{Three constraints that define a diffusion model}

A \textbf{Variational Diffusion Model} (VDM; Kingma et al.\ 2021;
Luo 2022) is an MHVAE with three additional constraints:

\begin{block}{Constraint 1 — Full-dimensional latents}
$\dim(\bx_t) = \dim(\bx_0)$ for all $t$.
No bottleneck; the latents live in \emph{data space}.
\end{block}

\begin{block}{Constraint 2 — Fixed Gaussian encoder (forward process)}
\[
q(\bx_t|\bx_{t-1})
= \Norm\!\bigl(\bx_t;\,\sqrt{\alpha_t}\,\bx_{t-1},\,(1-\alpha_t)\mathbf{I}\bigr),
\quad \alpha_t = 1-\beta_t\in(0,1).
\]
Not trained. The noise schedule $\{\beta_t\}$ is a design choice.
\end{block}

\begin{block}{Constraint 3 — Gaussian terminal distribution}
$\{\alpha_t\}$ is chosen so that
$q(\bx_T|\bx_0)\approx\Norm(\bm{0},\mathbf{I})$
as $t\to T$.
The model ``forgets'' the data completely.
\end{block}
\end{frame}

% =============================================================================
\section{The Forward Process}
% =============================================================================

% -----------------------------------------------------------------------------
\begin{frame}{Forward process: the Gaussian product rule}

\textbf{Key tool.} If $\bx = a\bm{u} + b\bm{v}$ with
$\bm{u}\sim\Norm(\bm{0},\mathbf{I})$ and $\bm{v}\sim\Norm(\bm{0},\mathbf{I})$
independent, then
\[
  \bx \sim \Norm\!\bigl(\bm{0},\,(a^2+b^2)\mathbf{I}\bigr).
\]

\textbf{Proof of closed-form marginal.}
Write one step as $\bx_t = \sqrt{\alpha_t}\,\bx_{t-1} + \sqrt{1-\alpha_t}\,\beps_{t-1}$
with $\beps_{t-1}\sim\Norm(\bm{0},\mathbf{I})$.  Unrolling two steps:
\begin{align*}
\bx_t
&= \sqrt{\alpha_t}\bigl(\sqrt{\alpha_{t-1}}\,\bx_{t-2}
   + \sqrt{1-\alpha_{t-1}}\,\beps_{t-2}\bigr)
   + \sqrt{1-\alpha_t}\,\beps_{t-1}\\
&= \sqrt{\alpha_t\alpha_{t-1}}\,\bx_{t-2}
   + \underbrace{\sqrt{\alpha_t(1-\alpha_{t-1})}\,\beps_{t-2}
     + \sqrt{1-\alpha_t}\,\beps_{t-1}}_{\sim\,\Norm(\bm{0},\,
       [\alpha_t(1-\alpha_{t-1})+(1-\alpha_t)]\mathbf{I})}\\
&= \sqrt{\alpha_t\alpha_{t-1}}\,\bx_{t-2}
   + \sqrt{1-\alpha_t\alpha_{t-1}}\;\bar{\beps}
\end{align*}
since $\alpha_t(1-\alpha_{t-1})+(1-\alpha_t) = 1-\alpha_t\alpha_{t-1}$.
\end{frame}

% -----------------------------------------------------------------------------
\begin{frame}{Closed-form forward marginal $q(\bx_t|\bx_0)$}

Define $\bar\alpha_t = \prod_{s=1}^{t}\alpha_s$.  Continuing the induction:

\begin{block}{Closed-form marginal (key result)}
\[
\boxed{
  q(\bx_t|\bx_0)
  = \Norm\!\bigl(\bx_t;\;\sqrt{\bar\alpha_t}\,\bx_0,\;
    (1-\bar\alpha_t)\mathbf{I}\bigr)
}
\]
Equivalently, we can sample $\bx_t$ \emph{directly} in one step:
\[
  \bx_t = \sqrt{\bar\alpha_t}\,\bx_0 + \sqrt{1-\bar\alpha_t}\,\beps,
  \qquad \beps\sim\Norm(\bm{0},\mathbf{I}).
\]
\end{block}

\medskip
\textbf{Interpretation.}
\begin{itemize}
  \item $\sqrt{\bar\alpha_t}$ is the \emph{signal coefficient}: decays toward 0.
  \item $\sqrt{1-\bar\alpha_t}$ is the \emph{noise coefficient}: grows toward 1.
  \item At $t=0$: $\bx_0 = \bx_0$ (clean).
  \item At $t=T$: $\bx_T\approx\beps$ (pure noise).
\end{itemize}
No forward-pass through a network needed — this is just arithmetic.
\end{frame}

% -----------------------------------------------------------------------------
\begin{frame}{Signal-to-noise ratio and noise schedules}

Define the \textbf{signal-to-noise ratio}:
\[
  \mathrm{SNR}(t) = \frac{\bar\alpha_t}{1-\bar\alpha_t}.
\]
$\mathrm{SNR}(0)\gg 1$ (clean data); $\mathrm{SNR}(T)\approx 0$ (pure noise).
The VDM objective can be rewritten entirely in terms of $\mathrm{SNR}(t)$,
and the loss weight of each term equals $\mathrm{SNR}(t)-\mathrm{SNR}(t+1)$
(Kingma et al.\ 2021).

\bigskip
\begin{columns}[T]
\column{0.5\textwidth}
\textbf{Linear schedule} (Ho et al.\ 2020):
\[
  \beta_t = \beta_{\min} + \tfrac{t-1}{T-1}(\beta_{\max}-\beta_{\min})
\]
$\beta_{\min}=10^{-4}$, $\beta_{\max}=0.02$, $T=1000$.
\smallskip

Near-zero SNR early, then rapid collapse.

\column{0.48\textwidth}
\textbf{Cosine schedule} (Nichol \& Dhariwal 2021):
\[
  \bar\alpha_t = \frac{f(t)}{f(0)},\quad
  f(t) = \cos^2\!\!\left(\frac{t/T + s}{1+s}\cdot\frac{\pi}{2}\right)
\]
$s=0.008$. SNR decreases more gradually, improving training at low $t$.
\end{columns}

\medskip
The cosine schedule avoids the very small $\beta_t$ values that waste
early training steps on near-identical images.
\end{frame}

% =============================================================================
\section{The Forward Posterior}
% =============================================================================

% -----------------------------------------------------------------------------
\begin{frame}{Tractable forward posterior $q(\bx_{t-1}|\bx_t,\bx_0)$}

The \emph{unconditional} reverse $q(\bx_{t-1}|\bx_t)$ is intractable.
But conditioned on $\bx_0$ it is tractable via Bayes' rule:
\[
  q(\bx_{t-1}|\bx_t,\bx_0)
  = \frac{q(\bx_t|\bx_{t-1},\bx_0)\,q(\bx_{t-1}|\bx_0)}{q(\bx_t|\bx_0)}
  = \frac{q(\bx_t|\bx_{t-1})\,q(\bx_{t-1}|\bx_0)}{q(\bx_t|\bx_0)}.
\]
The last equality uses the Markov property: $q(\bx_t|\bx_{t-1},\bx_0)=q(\bx_t|\bx_{t-1})$.

\medskip
Each factor is Gaussian:
\begin{align*}
q(\bx_t|\bx_{t-1})    &= \Norm(\sqrt{\alpha_t}\,\bx_{t-1},\,(1-\alpha_t)\mathbf{I})\\
q(\bx_{t-1}|\bx_0)   &= \Norm(\sqrt{\bar\alpha_{t-1}}\,\bx_0,\,(1-\bar\alpha_{t-1})\mathbf{I})\\
q(\bx_t|\bx_0)        &= \Norm(\sqrt{\bar\alpha_t}\,\bx_0,\,(1-\bar\alpha_t)\mathbf{I})
\end{align*}
Substituting and completing the square in $\bx_{t-1}$
yields a Gaussian in $\bx_{t-1}$.
\end{frame}

% -----------------------------------------------------------------------------
\begin{frame}{Forward posterior: completing the square}

The log of the numerator (as a function of $\bx_{t-1}$) is:
\begin{align*}
&-\frac{(\bx_t - \sqrt{\alpha_t}\,\bx_{t-1})^2}{2(1-\alpha_t)}
 -\frac{(\bx_{t-1} - \sqrt{\bar\alpha_{t-1}}\,\bx_0)^2}{2(1-\bar\alpha_{t-1})}\\[4pt]
=\;&-\frac{1}{2}\left[
  \frac{\alpha_t}{1-\alpha_t}\bx_{t-1}^2
  - \frac{2\sqrt{\alpha_t}\,\bx_t}{1-\alpha_t}\bx_{t-1}
  + \frac{1}{1-\bar\alpha_{t-1}}\bx_{t-1}^2
  - \frac{2\sqrt{\bar\alpha_{t-1}}\,\bx_0}{1-\bar\alpha_{t-1}}\bx_{t-1}
  + \text{const}\right]
\end{align*}

Grouping the $\bx_{t-1}^2$ and $\bx_{t-1}$ terms gives a precision:
\[
  \frac{1}{\tilde\beta_t}
  = \frac{\alpha_t}{1-\alpha_t} + \frac{1}{1-\bar\alpha_{t-1}}
  = \frac{\alpha_t(1-\bar\alpha_{t-1}) + (1-\alpha_t)}{(1-\alpha_t)(1-\bar\alpha_{t-1})}
  = \frac{1-\bar\alpha_t}{(1-\alpha_t)(1-\bar\alpha_{t-1})}.
\]
Therefore:
\[
  \tilde\beta_t
  = \frac{(1-\alpha_t)(1-\bar\alpha_{t-1})}{1-\bar\alpha_t}
  = \frac{1-\bar\alpha_{t-1}}{1-\bar\alpha_t}\,\beta_t.
\]
\end{frame}

% -----------------------------------------------------------------------------
\begin{frame}{Forward posterior: the mean $\tilde{\bmu}_t$}

The mean of $q(\bx_{t-1}|\bx_t,\bx_0)$ is read off from the linear term:
\[
  \frac{\tilde{\bmu}_t}{\tilde\beta_t}
  = \frac{\sqrt{\alpha_t}\,\bx_t}{1-\alpha_t}
  + \frac{\sqrt{\bar\alpha_{t-1}}\,\bx_0}{1-\bar\alpha_{t-1}}.
\]
Multiplying through by $\tilde\beta_t$:
\begin{align*}
\tilde{\bmu}_t(\bx_t,\bx_0)
&= \frac{\sqrt{\alpha_t}(1-\bar\alpha_{t-1})}{1-\bar\alpha_t}\,\bx_t
 + \frac{\sqrt{\bar\alpha_{t-1}}\,\beta_t}{1-\bar\alpha_t}\,\bx_0
\end{align*}

\begin{block}{Forward posterior (exact result)}
\[
\boxed{
q(\bx_{t-1}|\bx_t,\bx_0)
= \Norm\!\bigl(\bx_{t-1};\;\tilde{\bmu}_t(\bx_t,\bx_0),\;\tilde\beta_t\mathbf{I}\bigr)
}
\]
\vspace{-2mm}
\[
\tilde{\bmu}_t(\bx_t,\bx_0)
= \frac{\sqrt{\alpha_t}(1-\bar\alpha_{t-1})}{1-\bar\alpha_t}\,\bx_t
+ \frac{\sqrt{\bar\alpha_{t-1}}\,\beta_t}{1-\bar\alpha_t}\,\bx_0,
\qquad
\tilde\beta_t = \frac{1-\bar\alpha_{t-1}}{1-\bar\alpha_t}\,\beta_t
\]
\end{block}

This is the \emph{target} the learned reverse process must match.
\end{frame}

% =============================================================================
\section{The ELBO and Its Decomposition}
% =============================================================================

% -----------------------------------------------------------------------------
\begin{frame}{VDM ELBO: full decomposition}

Substituting the VDM structure into the MHVAE ELBO and applying
the Markov property at each level (Luo 2022, Sec.\ 3):

\begin{align*}
\log p(\bx_0)
&\geq
\underbrace{\E_{q(\bx_1|\bx_0)}
  \!\left[\log p_{\btheta}(\bx_0|\bx_1)\right]}_{\mathcal{L}_0:\;\text{reconstruction}}\\
&\quad
-\underbrace{\DKL\!\bigl(q(\bx_T|\bx_0)\,\|\,p(\bx_T)\bigr)}_{\mathcal{L}_T:\;\text{prior matching}}\\
&\quad
-\underbrace{\sum_{t=2}^{T}\E_{q(\bx_t|\bx_0)}\!
  \left[\DKL\!\bigl(q(\bx_{t-1}|\bx_t,\bx_0)\,\|\,
               p_{\btheta}(\bx_{t-1}|\bx_t)\bigr)\right]}_{\mathcal{L}_{1:T-1}:\;\text{denoising matching}}
\end{align*}

\begin{itemize}
  \item $\mathcal{L}_T = 0$: the noise schedule is designed so
        $q(\bx_T|\bx_0)\approx p(\bx_T)=\Norm(\bm{0},\mathbf{I})$. No parameters here.
  \item $\mathcal{L}_0$: a single reconstruction step; treated separately.
  \item $\mathcal{L}_{1:T-1}$: the dominant $T-1$ terms, each a Gaussian KL.
\end{itemize}
\end{frame}

% -----------------------------------------------------------------------------
\begin{frame}{Why condition on $\bx_0$ reduces variance}

\textbf{First attempt (naïve ELBO)} uses $q(\bx_{t-1}|\bx_t)$ in the KL,
which requires an expectation over \emph{two} random variables per term.
This has high Monte Carlo variance.

\medskip
\textbf{Second attempt (conditioned ELBO)} uses $q(\bx_{t-1}|\bx_t,\bx_0)$.
By the tower property:
\[
  \E_{q(\bx_{t-1},\bx_t|\bx_0)}\!\left[
    \DKL\!\bigl(q(\bx_{t-1}|\bx_t)\,\|\,p_{\btheta}(\bx_{t-1}|\bx_t)\bigr)\right]
  \;\geq\;
  \E_{q(\bx_t|\bx_0)}\!\left[
    \DKL\!\bigl(q(\bx_{t-1}|\bx_t,\bx_0)\,\|\,p_{\btheta}(\bx_{t-1}|\bx_t)\bigr)\right].
\]
The conditioned form is a tighter bound \emph{and} requires sampling only
$\bx_t\sim q(\bx_t|\bx_0)$, which is one-shot via the closed-form marginal.
Each KL term now involves only a \textbf{single Monte Carlo draw}.

\begin{alertblock}{Practical consequence}
We compute $\bx_t = \sqrt{\bar\alpha_t}\,\bx_0 + \sqrt{1-\bar\alpha_t}\,\beps$
and evaluate the KL analytically — no expensive nested sampling.
\end{alertblock}
\end{frame}

% -----------------------------------------------------------------------------
\begin{frame}{Gaussian KL: closed-form evaluation}

Each denoising term is a KL between two Gaussians.
For general Gaussians $\mathcal{N}(\bmu_1,\Sigma_1)$ and
$\mathcal{N}(\bmu_2,\Sigma_2)$ in $\R^d$:
\[
  \DKL\!\bigl(\Norm(\bmu_1,\Sigma_1)\,\|\,\Norm(\bmu_2,\Sigma_2)\bigr)
  = \half\!\left[
    \log\frac{|\Sigma_2|}{|\Sigma_1|}
    - d
    + \mathrm{tr}(\Sigma_2^{-1}\Sigma_1)
    + (\bmu_2-\bmu_1)^\top\Sigma_2^{-1}(\bmu_2-\bmu_1)
  \right].
\]

\medskip
In our case both distributions have variance $\tilde\beta_t\mathbf{I}$,
so the $\log|\cdot|$ and $\mathrm{tr}(\cdot)$ terms cancel exactly:
\[
  \DKL\!\bigl(q(\bx_{t-1}|\bx_t,\bx_0)\,\|\,p_{\btheta}(\bx_{t-1}|\bx_t)\bigr)
  = \frac{1}{2\tilde\beta_t}
    \bigl\|\tilde{\bmu}_t(\bx_t,\bx_0) - \bmu_{\btheta}(\bx_t,t)\bigr\|^2
  + C
\]
where $p_{\btheta}(\bx_{t-1}|\bx_t)=\Norm(\bmu_{\btheta}(\bx_t,t),\tilde\beta_t\mathbf{I})$
and $C$ does not depend on $\btheta$.

\medskip
\textbf{Training reduces to matching means.}
The network only needs to predict $\tilde{\bmu}_t$.
\end{frame}

% =============================================================================
\section{Noise Prediction and the Training Objective}
% =============================================================================

% -----------------------------------------------------------------------------
\begin{frame}{Reparameterising $\tilde{\bmu}_t$ in terms of noise}

From the closed-form forward marginal:
$\bx_t = \sqrt{\bar\alpha_t}\,\bx_0 + \sqrt{1-\bar\alpha_t}\,\beps$,
we can solve for $\bx_0$:
\[
  \bx_0 = \frac{1}{\sqrt{\bar\alpha_t}}\bigl(\bx_t - \sqrt{1-\bar\alpha_t}\,\beps\bigr).
\]
Substituting into $\tilde{\bmu}_t(\bx_t,\bx_0)$:
\begin{align*}
\tilde{\bmu}_t(\bx_t,\beps)
&= \frac{\sqrt{\alpha_t}(1-\bar\alpha_{t-1})}{1-\bar\alpha_t}\,\bx_t
 + \frac{\sqrt{\bar\alpha_{t-1}}\,\beta_t}{1-\bar\alpha_t}
   \cdot\frac{\bx_t - \sqrt{1-\bar\alpha_t}\,\beps}{\sqrt{\bar\alpha_t}}\\
&= \frac{\sqrt{\alpha_t}(1-\bar\alpha_{t-1}) + \beta_t\sqrt{\bar\alpha_{t-1}/\bar\alpha_t}}{1-\bar\alpha_t}\,\bx_t
 - \frac{\beta_t\sqrt{\bar\alpha_{t-1}/\bar\alpha_t}\,\sqrt{1-\bar\alpha_t}}{1-\bar\alpha_t}\,\beps
\end{align*}

Using $\bar\alpha_t = \bar\alpha_{t-1}\alpha_t$ so $\sqrt{\bar\alpha_{t-1}/\bar\alpha_t}=1/\sqrt{\alpha_t}$,
and collecting terms (see next slide):
\end{frame}

% -----------------------------------------------------------------------------
\begin{frame}{The noise-parameterised mean}

After collecting terms in $\bx_t$ and $\beps$:
\[
  \tilde{\bmu}_t
  = \frac{1}{\sqrt{\alpha_t}}\!\left(
      \bx_t - \frac{\beta_t}{\sqrt{1-\bar\alpha_t}}\,\beps
    \right).
\]

\textbf{Verification.}
Numerator of $\bx_t$ coefficient:
$\alpha_t(1-\bar\alpha_{t-1}) + \beta_t/\alpha_t \cdot \bar\alpha_{t-1}/\bar\alpha_{t-1}$...
simplifies to $(1-\bar\alpha_t)/\sqrt{\alpha_t}$, giving coefficient $1/\sqrt{\alpha_t}$.\quad$\checkmark$

\medskip
If the network predicts noise
$\beps_{\btheta}(\bx_t,t)\approx\beps$, then choose:
\[
\bmu_{\btheta}(\bx_t,t)
= \frac{1}{\sqrt{\alpha_t}}\!\left(
    \bx_t - \frac{\beta_t}{\sqrt{1-\bar\alpha_t}}\,\beps_{\btheta}(\bx_t,t)
  \right).
\]

The denoising KL becomes:
\[
  \mathcal{L}_{t-1}
  = \E_{\bx_0,\beps}\!\left[
    \frac{\beta_t^2}{2\tilde\beta_t\,\alpha_t(1-\bar\alpha_t)}\,
    \bigl\|\beps - \beps_{\btheta}(\bx_t,t)\bigr\|^2
  \right]
\]
where $\bx_t = \sqrt{\bar\alpha_t}\,\bx_0 + \sqrt{1-\bar\alpha_t}\,\beps$.
\end{frame}

% -----------------------------------------------------------------------------
\begin{frame}{The simplified DDPM training objective}

Ho et al.\ (2020) observed empirically that \textbf{dropping the time-dependent
weighting} $\beta_t^2/(2\tilde\beta_t\alpha_t(1-\bar\alpha_t))$ and averaging
uniformly over $t$ gives a simpler objective with \emph{better} sample quality:

\begin{block}{DDPM training loss $\mathcal{L}_{\mathrm{simple}}$}
\[
\boxed{
\mathcal{L}_{\mathrm{simple}}
= \E_{t\sim\mathrm{U}[1,T],\;\bx_0,\;\beps\sim\Norm(\bm{0},\mathbf{I})}\!\left[
  \bigl\|\beps
         - \beps_{\btheta}\!\left(
             \sqrt{\bar\alpha_t}\,\bx_0
             + \sqrt{1-\bar\alpha_t}\,\beps,\;t
           \right)\bigr\|^2
\right]
}
\]
\end{block}

\medskip
\textbf{Intuition.}
Dropping the weight up-weights early steps (small $t$, small noise, easy
to denoise) relative to later steps. This has a regularising effect:
the network focuses more on the hard tasks where its predictions matter.
The original weighted sum is still a valid ELBO; the simple loss is a
reweighted ELBO that happens to work better in practice.
\end{frame}

% -----------------------------------------------------------------------------
\begin{frame}{Three equivalent parameterisations}

Because $\bx_t,\bx_0,\beps$ are related by a linear equation,
predicting any one of them is equivalent:

\begin{center}
\renewcommand{\arraystretch}{1.35}
\begin{tabular}{llp{5.2cm}}
\toprule
\textbf{Predict} & \textbf{Objective} & \textbf{Notes}\\
\midrule
$\beps$ (noise) &
  $\|\beps - \beps_{\btheta}(\bx_t,t)\|^2$ &
  DDPM (Ho et al.\ 2020); most common\\[4pt]
$\bx_0$ (data) &
  $\|\bx_0 - \hat\bx_{\btheta}(\bx_t,t)\|^2$ &
  Allows explicit image prior; can over-smooth\\[4pt]
$\bm{v} = \sqrt{\bar\alpha_t}\,\beps - \sqrt{1-\bar\alpha_t}\,\bx_0$ &
  $\|\bm{v} - \bm{v}_{\btheta}(\bx_t,t)\|^2$ &
  Velocity (Salimans \& Ho 2022); stable near $t=0$\\
\bottomrule
\end{tabular}
\end{center}

\medskip
\textbf{Conversions at inference:}
\[
\hat\bx_0 = \frac{\bx_t - \sqrt{1-\bar\alpha_t}\,\beps_{\btheta}}{\sqrt{\bar\alpha_t}},
\qquad
\hat\beps = \frac{\bx_t - \sqrt{\bar\alpha_t}\,\hat\bx_0}{\sqrt{1-\bar\alpha_t}}.
\]
All three targets give the same denoising KL when substituted back.
\end{frame}

% =============================================================================
\section{Connection to Score Matching}
% =============================================================================

% -----------------------------------------------------------------------------
\begin{frame}{The score function}

The \textbf{score} of a distribution $p(\bx)$ is
$\nabla_{\bx}\log p(\bx)$.
It points towards regions of high probability density.

\medskip
\textbf{Score of the forward-diffused distribution.}
Since
$q(\bx_t|\bx_0) = \Norm(\sqrt{\bar\alpha_t}\,\bx_0,(1-\bar\alpha_t)\mathbf{I})$,
its log-density is:
\[
  \log q(\bx_t|\bx_0)
  = -\frac{1}{2(1-\bar\alpha_t)}\|\bx_t - \sqrt{\bar\alpha_t}\,\bx_0\|^2 + C.
\]
Differentiating with respect to $\bx_t$:
\[
  \nabla_{\bx_t}\log q(\bx_t|\bx_0)
  = -\frac{\bx_t - \sqrt{\bar\alpha_t}\,\bx_0}{1-\bar\alpha_t}
  = -\frac{\beps}{\sqrt{1-\bar\alpha_t}}.
\]
The score is proportional to the \emph{negative noise}.
\end{frame}

% -----------------------------------------------------------------------------
\begin{frame}{DDPM as denoising score matching}

Define the learned score network $\bm{s}_{\btheta}(\bx_t,t)$.
The \textbf{score matching objective} minimises:
\[
  \E_{t,\bx_t}\!\left[\bigl\|\nabla_{\bx_t}\log q(\bx_t) - \bm{s}_{\btheta}(\bx_t,t)\bigr\|^2\right].
\]

The \textbf{denoising score matching} (Vincent 2011) substitution:
\[
  \nabla_{\bx_t}\log q(\bx_t)
  = \E_{q(\bx_0|\bx_t)}\!\left[\nabla_{\bx_t}\log q(\bx_t|\bx_0)\right]
  = -\E_{q(\bx_0|\bx_t)}\!\left[\frac{\beps}{\sqrt{1-\bar\alpha_t}}\right].
\]

With the optimal score network $\bm{s}^*_{\btheta}(\bx_t,t)=-\beps^*/\sqrt{1-\bar\alpha_t}$:
\[
  \bm{s}_{\btheta}(\bx_t,t)
  = -\frac{\beps_{\btheta}(\bx_t,t)}{\sqrt{1-\bar\alpha_t}}
  \;\;\Longrightarrow\;\;
  \bm{s}^*_{\btheta}\approx\nabla_{\bx_t}\log q(\bx_t).
\]

\begin{alertblock}{Equivalence}
\textbf{DDPM noise prediction = denoising score matching.}
The DDPM network implicitly learns the score of the noisy data distribution
at every noise level simultaneously.
\end{alertblock}
\end{frame}

% -----------------------------------------------------------------------------
\begin{frame}{From score to SDE: the continuous-time limit}

Song et al.\ (2021) embed the discrete chain into a continuous SDE.
As $T\to\infty$ with $\beta_t\to\beta(t)\,dt$, the forward process becomes:
\[
  d\bx = -\tfrac{1}{2}\beta(t)\,\bx\,dt + \sqrt{\beta(t)}\,d\bm{W}_t,
\]
where $\bm{W}_t$ is Brownian motion.  By Anderson (1982), the
\textbf{reverse SDE} is:
\[
  d\bx = \bigl[-\tfrac{1}{2}\beta(t)\,\bx
               - \beta(t)\,\nabla_{\bx}\log p_t(\bx)\bigr]\,dt
         + \sqrt{\beta(t)}\,d\bar{\bm{W}}_t
\]
where $\bar{\bm{W}}_t$ is reverse-time Brownian motion.
Replacing the score by the learned network gives:
\[
d\bx = \bigl[-\tfrac{1}{2}\beta(t)\,\bx
             + \beta(t)\,\tfrac{\beps_{\btheta}(\bx,t)}{\sqrt{1-\bar\alpha_t}}\bigr]\,dt
       + \sqrt{\beta(t)}\,d\bar{\bm{W}}_t.
\]
The same network $\beps_{\btheta}$ that DDPM trains is the score estimator
that drives the reverse SDE.
\end{frame}

% =============================================================================
\section{The Reverse Process and Sampling}
% =============================================================================

% -----------------------------------------------------------------------------
\begin{frame}{The learned reverse process}

Choose the learned reverse as:
\[
  p_{\btheta}(\bx_{t-1}|\bx_t)
  = \Norm\!\bigl(\bx_{t-1};\;\bmu_{\btheta}(\bx_t,t),\;\tilde\beta_t\mathbf{I}\bigr)
\]
with $\bmu_{\btheta}$ the noise-parameterised mean from Section~4.

\medskip
\textbf{Why use $\tilde\beta_t$ as the reverse variance?}
The true posterior variance $\tilde\beta_t$ lies between
$\beta_t$ (upper bound, for $q(\bx_{t-1}|\bx_t)$) and
$\tilde\beta_t$ (exact, for $q(\bx_{t-1}|\bx_t,\bx_0)$).
Ho et al.\ fixed it to $\tilde\beta_t$; later work (Nichol \& Dhariwal 2021)
learns it as an interpolation:
\[
  \Sigma_{\btheta}(\bx_t,t)
  = \exp\bigl(v\log\beta_t + (1-v)\log\tilde\beta_t\bigr),
  \quad v = v_{\btheta}(\bx_t,t)\in[0,1].
\]

\begin{block}{DDPM reverse step}
\[
  \bx_{t-1} = \frac{1}{\sqrt{\alpha_t}}\!\left(
    \bx_t - \frac{\beta_t}{\sqrt{1-\bar\alpha_t}}\,\beps_{\btheta}(\bx_t,t)
  \right)
  + \sqrt{\tilde\beta_t}\,\bz,\quad
  \bz\sim\Norm(\bm{0},\mathbf{I})\;(t>1),\;\bz=\bm{0}\;(t=1).
\]
\end{block}
\end{frame}

% -----------------------------------------------------------------------------
\begin{frame}{DDPM reverse algorithm}

\begin{block}{Complete DDPM sampling procedure}
\begin{enumerate}
  \setlength\itemsep{5pt}
  \item Sample $\bx_T\sim\Norm(\bm{0},\mathbf{I})$.
  \item For $t = T,\, T-1,\,\ldots,\,1$:
    \begin{enumerate}
      \item $\bz\sim\Norm(\bm{0},\mathbf{I})$ if $t>1$;\quad $\bz=\bm{0}$ if $t=1$.
      \item Predict noise: $\hat\beps = \beps_{\btheta}(\bx_t,t)$.
      \item Update:
        $\displaystyle\bx_{t-1}
         = \frac{1}{\sqrt{\alpha_t}}\!\left(
             \bx_t - \frac{\beta_t}{\sqrt{1-\bar\alpha_t}}\,\hat\beps
           \right)
           + \sqrt{\tilde\beta_t}\,\bz$.
    \end{enumerate}
  \item Return $\bx_0$.
\end{enumerate}
\end{block}

\medskip
\begin{itemize}
  \item Each step is one network evaluation plus a cheap Gaussian update.
  \item With $T=1000$: slow but high quality.
  \item \textbf{No discriminator}, no adversarial dynamics, no mode collapse.
\end{itemize}
\end{frame}

% =============================================================================
\section{Faster Sampling: DDIM}
% =============================================================================

% -----------------------------------------------------------------------------
\begin{frame}{DDIM: deterministic ODE sampling}

Song et al.\ (2020) observe that the DDPM objective trains a score network
that is \emph{independent of the sampling procedure}. One can therefore
replace the stochastic reverse SDE with a deterministic ODE:

\[
  d\bx = \left[\frac{\bx - \sqrt{1-\bar\alpha}\,\beps_{\btheta}(\bx,\bar\alpha)}
               {\sqrt{\bar\alpha}}\right] d\sqrt{\bar\alpha}
\]

Discretising over a \emph{subsequence} $\tau_1<\tau_2<\cdots<\tau_S$ of
timesteps (with $S\ll T$):

\begin{block}{DDIM update rule}
\[
  \bx_{\tau_{i-1}}
  = \sqrt{\bar\alpha_{\tau_{i-1}}}\,\hat\bx_0(\bx_{\tau_i})
  + \sqrt{1-\bar\alpha_{\tau_{i-1}} - \sigma_{\tau_i}^2}\;\beps_{\btheta}(\bx_{\tau_i},\tau_i)
  + \sigma_{\tau_i}\,\bz
\]
where $\hat\bx_0(\bx_{\tau_i})
= \frac{\bx_{\tau_i}-\sqrt{1-\bar\alpha_{\tau_i}}\,\beps_{\btheta}(\bx_{\tau_i},\tau_i)}{\sqrt{\bar\alpha_{\tau_i}}}$
and $\sigma_{\tau_i} = \eta\sqrt{\tilde\beta_{\tau_i}}$.
\end{block}

\begin{itemize}
  \item $\eta=0$: fully deterministic. Same $\bx_T\to$same~$\bx_0$.
  \item $\eta=1$: recovers DDPM.
  \item $S=50$ gives quality comparable to $T=1000$ DDPM.
\end{itemize}
\end{frame}

% -----------------------------------------------------------------------------
\begin{frame}{DDIM: derivation of the update rule}

Start from the DDIM non-Markovian inference distribution (Song et al.\ 2020):
\[
  q_\sigma(\bx_{\tau_{i-1}}|\bx_{\tau_i},\bx_0)
  = \Norm\!\bigl(\bx_{\tau_{i-1}};\;
    \sqrt{\bar\alpha_{\tau_{i-1}}}\,\bx_0
    + \sqrt{1-\bar\alpha_{\tau_{i-1}}-\sigma^2}\;\tfrac{\bx_{\tau_i}-\sqrt{\bar\alpha_{\tau_i}}\,\bx_0}{\sqrt{1-\bar\alpha_{\tau_i}}},\;
    \sigma^2\mathbf{I}\bigr).
\]

Replacing $\bx_0$ by the predicted $\hat\bx_0$ and the exact noise direction
by $\beps_{\btheta}(\bx_{\tau_i},\tau_i)$ recovers the update rule on the
previous slide.

\medskip
\textbf{Why does this work?}
The DDIM update is a first-order ODE solver for the probability-flow ODE.
The ODE and the reverse SDE share the same marginals
$p_t(\bx_t)$ at every time $t$.
Since $\beps_{\btheta}$ approximates the score of $p_t$,
integrating the ODE gives the same distribution over $\bx_0$
as ancestral sampling from the reverse SDE — but in far fewer steps.
\end{frame}

% =============================================================================
\section{Latent Diffusion and Connections}
% =============================================================================

% -----------------------------------------------------------------------------
\begin{frame}{Latent Diffusion Models (Stable Diffusion)}

\textbf{Problem.} Running DDPM in pixel space for high-resolution images
($512\times512\times3 = 786\,432$ dimensions) is prohibitively expensive.

\medskip
\textbf{Solution} (Rombach et al.\ 2022):

\begin{enumerate}
  \setlength\itemsep{4pt}
  \item Train a \textbf{VAE}: compress $\bx\in\R^{H\times W\times 3}$
        to $\bz\in\R^{h\times w\times c}$ with $h=H/f$, $w=W/f$ ($f=8$),
        $c=4$. Compression factor $\approx 48\times$.
  \item Run \textbf{DDPM in latent space} $\bz$: much cheaper.
  \item Decode with the VAE: $\hat\bz_0\to\hat\bx$.
\end{enumerate}

\begin{block}{Why this works}
\begin{itemize}
  \item The VAE provides a perceptually smooth, semantic latent space.
  \item Diffusion in $\bz$ avoids modelling low-level pixel correlations.
  \item Conditioning (text, depth map, class) enters via cross-attention in
        the U-Net denoiser operating on $\bz$.
  \item The VAE's KL regularisation ensures $\bz$ is smooth enough for a
        diffusion model to learn a clean manifold.
\end{itemize}
\end{block}
\end{frame}

% -----------------------------------------------------------------------------
\begin{frame}{VAE vs.\ Diffusion: mathematical parallel}

Both maximise the same variational lower bound on $\log p(\bx)$:
\[
  \log p(\bx) \geq \ELBO(\bphi,\btheta)
  = \underbrace{\E[\log p_{\btheta}(\bx|\bz)]}_{\text{reconstruction}}
  - \underbrace{\DKL(q_{\bphi}\,\|\,p)}_{\text{regularisation}}.
\]
Their differences arise from how they apply this framework:

\begin{center}
\renewcommand{\arraystretch}{1.3}
\begin{tabular}{lll}
\toprule
\textbf{Design axis} & \textbf{VAE} & \textbf{DDPM}\\
\midrule
Latent levels         & 1            & $T = 300$--$1000$\\
Encoder learned?      & Yes          & No (fixed schedule)\\
Bottleneck?           & Yes ($d_h\ll d_x$) & No ($d_h = d_x$)\\
What does network predict? & $\bx$ from $\bh$ & $\beps$ from $(\bx_t,t)$\\
KL terms              & 1            & $T-1$\\
Training loss         & BCE + $\DKL$ & MSE (noise)\\
Sampling cost         & 1 decoder call & $T$ denoiser calls\\
\bottomrule
\end{tabular}
\end{center}

\begin{block}{Key insight}
DDPM = VAE with (a) $T$ latent levels, (b) fixed encoder, (c) no bottleneck.
\end{block}
\end{frame}

% -----------------------------------------------------------------------------
\begin{frame}{VAE vs.\ Diffusion: pros and cons}

\begin{columns}[T]
\column{0.5\textwidth}
\begin{block}{VAE}
\textcolor{green!60!black}{\textbf{Pros:}}
\begin{itemize}\small
  \item Fast inference and generation (1 pass)
  \item Compact structured latent space
  \item Stable training with closed-form KL
  \item Good for downstream tasks (clustering, interpolation)
  \item Scalable to large datasets
\end{itemize}
\textcolor{red!60!black}{\textbf{Cons:}}
\begin{itemize}\small
  \item \textbf{Blurry samples} ($\ell_2$-averaging)
  \item Posterior collapse risk
  \item ELBO gap — log-likelihood never fully recovered
  \item Mean-field approximation limits expressivity
\end{itemize}
\end{block}

\column{0.5\textwidth}
\begin{block}{Diffusion (DDPM)}
\textcolor{green!60!black}{\textbf{Pros:}}
\begin{itemize}\small
  \item \textbf{State-of-the-art} sample quality
  \item No mode collapse — full distribution coverage
  \item Stable training (MSE objective, no adversary)
  \item Theoretically grounded via ELBO + score matching
  \item Flexible conditioning
\end{itemize}
\textcolor{red!60!black}{\textbf{Cons:}}
\begin{itemize}\small
  \item \textbf{Slow} sampling ($T$ network evaluations)
  \item No compact latent for downstream tasks
  \item High memory for large~$T$
  \item Less natural for representation learning
\end{itemize}
\end{block}
\end{columns}
\end{frame}

% =============================================================================
\section{Summary and Exercises}
% =============================================================================

% -----------------------------------------------------------------------------
\begin{frame}{Key equations: complete reference}

\begin{center}\small
\renewcommand{\arraystretch}{1.4}
\begin{tabular}{ll}
\toprule
\textbf{Concept} & \textbf{Formula}\\
\midrule
Forward step &
  $q(\bx_t|\bx_{t-1}) = \Norm(\sqrt{\alpha_t}\,\bx_{t-1},\,(1-\alpha_t)\mathbf{I})$\\
Forward marginal &
  $q(\bx_t|\bx_0) = \Norm(\sqrt{\bar\alpha_t}\,\bx_0,\,(1-\bar\alpha_t)\mathbf{I})$\\
One-shot sample &
  $\bx_t = \sqrt{\bar\alpha_t}\,\bx_0 + \sqrt{1-\bar\alpha_t}\,\beps$\\
Forward posterior mean &
  $\tilde\bmu_t = \frac{\sqrt{\alpha_t}(1-\bar\alpha_{t-1})}{1-\bar\alpha_t}\bx_t
    + \frac{\sqrt{\bar\alpha_{t-1}}\beta_t}{1-\bar\alpha_t}\bx_0$\\
Forward posterior var. &
  $\tilde\beta_t = \frac{1-\bar\alpha_{t-1}}{1-\bar\alpha_t}\,\beta_t$\\
Noise-param.\ mean &
  $\tilde\bmu_t = \frac{1}{\sqrt{\alpha_t}}\!\bigl(\bx_t - \frac{\beta_t}{\sqrt{1-\bar\alpha_t}}\beps\bigr)$\\
Training loss &
  $\mathcal{L}_{\mathrm{simple}} = \E\bigl[\|\beps - \beps_{\btheta}(\bx_t,t)\|^2\bigr]$\\
Score relation &
  $\bm{s}_{\btheta}(\bx_t,t) = -\beps_{\btheta}(\bx_t,t)/\sqrt{1-\bar\alpha_t}$\\
Reverse step &
  $\bx_{t-1} = \frac{1}{\sqrt{\alpha_t}}\!\bigl(\bx_t - \frac{\beta_t}{\sqrt{1-\bar\alpha_t}}\beps_{\btheta}\bigr) + \sqrt{\tilde\beta_t}\,\bz$\\
SNR &
  $\mathrm{SNR}(t) = \bar\alpha_t/(1-\bar\alpha_t)$\\
\bottomrule
\end{tabular}
\end{center}
\end{frame}

% -----------------------------------------------------------------------------
\begin{frame}{Exercises}

\begin{block}{Mathematical derivations}
\begin{enumerate}
  \setlength\itemsep{4pt}
  \item \textbf{Forward marginal.}
        Prove $q(\bx_t|\bx_0)=\Norm(\sqrt{\bar\alpha_t}\bx_0,(1-\bar\alpha_t)\mathbf{I})$
        by induction. Identify the variance at step $s+1$ from the variance at step $s$.

  \item \textbf{Forward posterior.}
        Apply Bayes' rule to derive $q(\bx_{t-1}|\bx_t,\bx_0)$.
        Show explicitly that completing the square yields
        mean $\tilde\bmu_t$ and variance $\tilde\beta_t$ as given.

  \item \textbf{Gaussian KL.}
        Derive $\DKL(\Norm(\bmu_1,\sigma^2\mathbf{I})\|\Norm(\bmu_2,\sigma^2\mathbf{I}))
        = \|\bmu_1-\bmu_2\|^2/(2\sigma^2)$ from the general KL formula,
        and verify that the log-determinant and trace terms cancel.

  \item \textbf{Noise reparameterisation.}
        Starting from $\tilde\bmu_t(\bx_t,\bx_0)$, substitute
        $\bx_0 = (\bx_t - \sqrt{1-\bar\alpha_t}\,\beps)/\sqrt{\bar\alpha_t}$
        and verify that the result equals
        $\frac{1}{\sqrt{\alpha_t}}(\bx_t - \frac{\beta_t}{\sqrt{1-\bar\alpha_t}}\beps)$.

  \item \textbf{Score matching.}
        Show that minimising the DDPM noise-prediction loss is equivalent
        to denoising score matching by writing the score of $q(\bx_t|\bx_0)$
        in terms of $\beps$ and identifying the score network.

  \item \textbf{ELBO decomposition.}
        Show that the MHVAE ELBO decomposes into
        $\mathcal{L}_0 + \mathcal{L}_T + \sum_{t=2}^T\mathcal{L}_{t-1}$
        by writing out the joint log-ratio and applying the chain rule of KL.
        Explain why $\mathcal{L}_T = 0$ by design.
\end{enumerate}
\end{block}
\end{frame}

% -----------------------------------------------------------------------------
\begin{frame}{References}

\begin{enumerate}\small
  \setlength\itemsep{4pt}
  \item D.\ Kingma \& M.\ Welling, \emph{Auto-Encoding Variational Bayes},
        ICLR 2014. \href{https://arxiv.org/abs/1312.6114}{arXiv:1312.6114}

  \item J.\ Sohl-Dickstein et al., \emph{Deep Unsupervised Learning using
        Nonequilibrium Thermodynamics}, ICML 2015.
        \href{https://arxiv.org/abs/1503.03585}{arXiv:1503.03585}

  \item J.\ Ho, A.\ Jain \& P.\ Abbeel, \emph{Denoising Diffusion Probabilistic
        Models}, NeurIPS 2020. \href{https://arxiv.org/abs/2006.11239}{arXiv:2006.11239}

  \item J.\ Song, C.\ Meng \& S.\ Ermon, \emph{Denoising Diffusion Implicit
        Models}, ICLR 2021. \href{https://arxiv.org/abs/2010.02502}{arXiv:2010.02502}

  \item Y.\ Song et al., \emph{Score-Based Generative Modeling through
        Stochastic Differential Equations}, ICLR 2021.
        \href{https://arxiv.org/abs/2011.13456}{arXiv:2011.13456}

  \item D.\ Kingma et al., \emph{Variational Diffusion Models}, NeurIPS 2021.
        \href{https://arxiv.org/abs/2107.00630}{arXiv:2107.00630}

  \item A.\ Nichol \& P.\ Dhariwal, \emph{Improved Denoising Diffusion
        Probabilistic Models}, ICML 2021.
        \href{https://arxiv.org/abs/2102.09672}{arXiv:2102.09672}

  \item R.\ Rombach et al., \emph{High-Resolution Image Synthesis with
        Latent Diffusion Models}, CVPR 2022.
        \href{https://arxiv.org/abs/2112.10752}{arXiv:2112.10752}

  \item C.\ Luo, \emph{Understanding Diffusion Models: A Unified Perspective},
        2022. \href{https://arxiv.org/abs/2208.11970}{arXiv:2208.11970}
\end{enumerate}
\end{frame}

\end{document}
